\documentclass{openjournal}

\usepackage[T1]{fontenc}
\usepackage{graphicx}
\usepackage{amsmath,amssymb}
\usepackage{hyperref}
\usepackage{tabularx}

\shorttitle{A Volume-Limited ALMA Search for Technosignatures within 20\,pc}
\shortauthors{White \& Dey}

\begin{document}

\title{A Volume-Limited Search for Technosignatures within 20\,pc Using the Public ALMA Archive}

\author{Glenn J. White\altaffilmark{1,2} and Robin Dey\altaffilmark{3}}
\altaffiltext{1}{School of Physical Sciences, The Open University, Walton Hall, Milton Keynes, MK7 6AA, UK; glenn.white@open.ac.uk}
\altaffiltext{2}{RAL Space, STFC Rutherford Appleton Laboratory, Harwell Campus, Didcot, OX11 0QX, UK}
\altaffiltext{3}{VBRL Holdings Inc.}

\begin{abstract}
We present the motivation, design, methodology and first results of a
volume-limited search for radio technosignatures among all stars within
20\,pc of the Sun that have public observations in the Atacama Large
Millimeter/submillimeter Array (ALMA) science archive. Unlike the great
majority of technosignature surveys conducted to date, which are either
wide, shallow all-sky surveys or targeted campaigns preferentially
selecting exoplanet hosts, young stellar objects, or other astrophysically
``interesting'' systems, our sample of 120 stars is every star within
20\,pc with public ALMA archival coverage -- selected on the presence of
ALMA data alone, irrespective of spectral type, age, multiplicity, or
known planetary or circumstellar-disc status, rather than on any of those
astrophysical properties. This sample is therefore unbiased with respect
to spectral type and disc/planet-host status \emph{within} the population
ALMA happens to have observed, though it is not, and does not claim to
be, a complete census of the local stellar neighbourhood: it represents
only $\sim$5\% of the $\sim$2400 individually catalogued stars within
20\,pc (Section~\ref{sec:sample}). We mine the ALMA archive
opportunistically: rather than requesting dedicated observing time, we
reprocess existing calibrated (and, where necessary, recalibrated)
visibility data to search simultaneously for narrowband Doppler-drifting
technosignature candidates and for anomalous broad-band continuum
emission, at frequencies ($\sim$\,90--950\,GHz) that remain almost
entirely unexplored by the technosignature literature. We supplement the
classical narrowband/continuum search with two further analyses that
exploit properties unique to this survey and unavailable to any
single-dish technosignature programme: a closure-phase point-source
vetting test for any candidate that crosses threshold, and a
population-level, spectral-type-normalised continuum anomaly check made
possible by the sample's volume-limited, unbiased design. As of this
version, 19 of the planned 120 target/bands have been searched to
completion. No credible technosignature has been found -- one narrowband
candidate initially flagged by the automated pipeline is identified as
the CO(1--0) transition and is not of technosignature origin -- and we
report EIRP upper limits in the range $1.6\times10^{13}$--$2.8\times10^{15}$\,W
(median $2.1\times10^{14}$\,W) and five continuum detections (one an
unresolved close binary), alongside fourteen deep non-detections
reaching a median 5$\sigma$ continuum limit of $\sim$\,0.6\,mJy. The
survey is ongoing; this paper will be updated as further targets are
completed.
\end{abstract}

\keywords{extraterrestrial intelligence --- radio lines: stars --- submillimetre: stars --- surveys --- astrobiology}

\section{Introduction}
\label{sec:intro}

The search for extraterrestrial intelligence (SETI) rests on the premise that a
technological civilisation, whether deliberately signalling or simply going about
its own affairs, may leave a detectable imprint on the electromagnetic spectrum: a
\emph{technosignature} \citep{Sheikh2020,SETIreview2025}. Since the foundational
proposal of \citet{CocconiMorrison1959} that microwave frequencies near the 21\,cm
hydrogen line would be an obvious ``meeting place'' for interstellar communication,
and Frank Drake's first observational test of the idea with Project Ozma
\citep{Drake1960}, the radio band has remained the dominant search space for
technosignatures, for reasons that are both historical and physically well
motivated: radio photons are cheap to produce in enormous numbers, propagate
through the interstellar medium and planetary atmospheres with little attenuation,
and (for a sufficiently narrowband signal) can in principle be distinguished from
any known natural astrophysical process by their spectral compactness and Doppler
dynamics \citep{Sheikh2020}.

Six decades on, radio technosignature searches have grown from single-target,
single-frequency pointings into surveys of many thousands of stars across broad
swathes of the radio spectrum. Yet, as \citet{Wright2018} demonstrated
quantitatively with their ``Cosmic Haystack'' formalism, even the combined efforts
of every major SETI programme to date have sampled only an infinitesimal fraction
of the plausible multi-dimensional parameter space of transmitter frequency,
bandwidth, power, location, and duty cycle. Persistent, non-trivial gaps remain:
almost the entire millimetre and submillimetre regime above $\sim$\,50\,GHz is
essentially unsearched; most surveys have targeted specific classes of star
(nearby exoplanet hosts, young stellar objects with protoplanetary discs, stars in
the ``Earth Transit Zone'') rather than an unbiased, complete sample of the local
stellar population; and only a small number of studies have systematically
reprocessed \emph{existing} interferometric archives for technosignatures, as
opposed to requesting dedicated telescope time.

This paper describes a new survey designed to address several of these gaps
simultaneously. We construct a volume-limited sample of every star within
20\,pc of the Sun with public ALMA archival coverage, irrespective of
spectral type, disc-hosting status, or known planets, and mine that
archival data for both narrowband technosignature candidates and
anomalous continuum emission, applying corrections for stellar proper
motion and for the range of Doppler drift rates plausible for a
transmitter on a planet in orbit around its host star. We further
supplement this classical search with two analyses that exploit
properties specific to this survey and unavailable to prior
technosignature work: interferometric closure-phase vetting of any
candidate signal, and a population-level anomaly check enabled by the
sample's unbiased, volume-limited design.
Section~\ref{sec:background} reviews the relevant
literature: the history of radio SETI (\S\ref{sec:history}), the current
generation of large-scale technosignature surveys (\S\ref{sec:modern}), the very
limited body of work at millimetre/submillimetre frequencies
(\S\ref{sec:mmsubmm}), the growing use of machine learning and anomaly detection
in this field (\S\ref{sec:ml}), the statistical frameworks used to interpret
survey non-detections (\S\ref{sec:stats}), the case for a volume-limited,
representative sample (\S\ref{sec:volumelimited}), and the prospects offered by
next-generation facilities (\S\ref{sec:future}). Section~\ref{sec:sample}
summarises our sample selection, Section~\ref{sec:method} our data-reduction and
search methodology (including the two new analyses),
Section~\ref{sec:results} presents our results to date, and
Section~\ref{sec:discussion} discusses their interpretation and the survey's
prospects. Full per-target tables are deferred to Appendix~\ref{app:table}
so as not to interrupt the main narrative with tabular material.

\section{Background}
\label{sec:background}

\subsection{Six decades of radio SETI: a brief history}
\label{sec:history}

Project Ozma \citep{Drake1960} observed two nearby Sun-like stars, $\tau$~Ceti and
$\epsilon$~Eridani, with the 26\,m Tatel telescope at Green Bank at 1420\,MHz for
a total of some 150 hours, and found no signal, but established both the
methodology and the symbolic starting point for the modern field. Through the
1960s--1980s a steady stream of targeted and sky-survey searches followed,
including Ozma~II \citep{Drake1986}, the Ohio State ``Big Ear'' survey (source of
the famous ``Wow!'' signal), and the META and BETA all-sky surveys conducted from
Harvard/Smithsonian facilities. Project Phoenix \citep{Backus2002}, run by the
SETI Institute from 1995--2004, was the most sensitive targeted survey of its era,
observing on the order of 800 nearby Sun-like stars using the Arecibo, Parkes and
Green Bank telescopes, and pioneered a rigorous methodology for distinguishing
candidate signals from terrestrial radio-frequency interference (RFI) through
simultaneous multi-site or multi-antenna verification. The Allen Telescope Array
\citep[ATA;][]{Tarter2001}, jointly developed by the SETI Institute and UC
Berkeley and first commissioned in 2007, was the first radio telescope
purpose-built from the outset for commensal, wide-band SETI observing, and remains
in active use, including in recent technosignature searches of interstellar
objects \citep{ATAISO2025} and TRAPPIST-1 \citep{ATATRAPPIST2024}. The
distributed-computing project SETI@home \citep{SETIhome2025} processed
public-participant-donated CPU cycles to search Arecibo sky-survey data for
narrowband and pulsed signals for two decades before its 2020 shutdown, and its
final data-analysis summary was only published in 2025.

Two complementary observing strategies have been used throughout this history:
\emph{sky surveys}, covering a large fraction of the celestial sphere at modest
sensitivity per pointing, and \emph{targeted searches}, concentrating integration
time on a preselected list of promising stars \citep[e.g.][]{Drake1974}. Both
strategies persist in the modern era, but -- as we discuss in
\S\ref{sec:volumelimited} -- the majority of targeted programmes have selected
their samples on some presumed indicator of habitability or interest (proximity,
Sun-like spectral type, known planets, or membership of the Earth Transit Zone),
rather than aiming for a complete, unbiased census of the local stellar
population.

\subsection{Modern large-scale radio technosignature surveys}
\label{sec:modern}

The current generation of technosignature science is dominated by the
Breakthrough Listen initiative, launched in 2015 with a ten-year, \$100\,million
commitment to survey the one million nearest stars, the full Galactic plane and
centre, the 100 nearest galaxies, and a sample of exoplanet and planet-candidate
hosts, primarily using the Robert C.\ Byrd Green Bank Telescope (GBT) and the
CSIRO Parkes ``Murriyang'' telescope, and now increasingly via commensal
observing on MeerKAT \citep{BLoverview}. Its first major published dataset,
\citet{Enriquez2017}, searched 692 nearby stars over 1.1--1.9\,GHz with the GBT
and set limits on narrowband transmitters of EIRP\,$>10^{13}$\,W; this was
followed by \citet{Price2020} (1327 stars, 1.10--3.45\,GHz),
\citet{Gajjar2019EarthTransitZone}-style searches of the restricted Earth Transit
Zone at 3.95--8.00\,GHz, and, most recently, a substantial expansion of automated,
commensal technosignature searching with MeerKAT that has processed some 1.5
million coherent beams between 2023 and 2026 \citep{BLMeerKAT2026}. In parallel,
\citet{Margot2023} conducted a large citizen-science-style search of 11\,680
stars with the GBT at 1.15--1.73\,GHz, explicitly developing a statistical
formalism for interpreting non-detections in terms of the fraction of stars that
could host a detectable transmitter, \emph{provided the observed sample is
representative of the underlying stellar population of the search volume} -- a
condition that, as we discuss below, is often not satisfied by targeted,
interest-selected samples.

Beyond Breakthrough Listen, the last five years have seen international radio
facilities extend technosignature searches to new frequency regimes and
observational strategies: the Five-hundred-metre Aperture Spherical Telescope
\citep[FAST;][]{Zhang2020, FAST3IATLAS2025} in China has conducted commensal
searches with unmatched raw sensitivity at L-band, including recent narrowband
searches of interstellar objects; LOFAR and the pathfinder NenuFAR have pushed
technosignature searches down to 10--190\,MHz, with LOFAR alone having surveyed
over 1.6 million stars in commensal mode and dual-site LOFAR configurations
enabling near-real-time RFI rejection \citep{LOFARdualsite2023}; the Sardinia
Radio Telescope has conducted one of the first dedicated high-frequency
($>$\,18\,GHz) surveys of nearby stars and the Galactic Centre
\citep{SRT2025}; and commensal technosignature back-ends are now operating or in
development at the Karl G.\ Jansky Very Large Array \citep[the COSMIC
system;][]{Tremblay2024}, MeerKAT \citep{Czech2021}, and the Murchison Widefield
Array \citep{Morrison2023}. Occultation-based strategies -- observing a star
during a planet-planet or planet-star occultation, when geometric alignment
transiently favours detection of leakage or deliberate signals -- have also grown
in popularity, exemplified by recent Breakthrough Listen searches of TESS-selected
eclipsing systems \citep{TESSeclipsing2025}.

\subsection{Millimetre and submillimetre technosignature searches}
\label{sec:mmsubmm}

In stark contrast to the radio regime below $\sim$\,10\,GHz, the millimetre and
submillimetre bands remain almost entirely unexplored for technosignatures.
\citet{Mason2024} conducted, to our knowledge, the first technosignature search
ever performed with ALMA, searching archival Band\,3 data (spectral windows
centred at 90.642 and 93.151\,GHz) for narrowband signals towards 28 Galactic
stars selected from Gaia~DR3. Their approach used a ``stellar bycatch'' method:
rather than targeting stars directly, they searched for serendipitous stars
falling within the undistorted field of view of four bright ALMA phase
calibrators, since these calibrators are observed frequently and are not
generally themselves of astrophysical interest. For the nearest star in their
sample they placed a limit of EIRP$_{\rm min} > 7\times10^{17}$\,W, and discussed
in detail the specific challenges of high-frequency technosignature work:
substantially larger Doppler drift rates at fixed line-of-sight acceleration
(since drift rate scales linearly with observing frequency), the resulting
sensitivity loss (``smearing'') for a fixed spectral and temporal resolution
unless a drift-rate search is performed, and an increased incidence of spectral
confusion from the dense molecular line forest present in protoplanetary discs
and star-forming regions at these frequencies. \citet{ATATRAPPIST2024} extended
technosignature work on the TRAPPIST-1 system to the Allen Telescope Array,
underlining the continued observational interest in this particular target but
also its status as one of very few systems that has now been searched for
technosignatures across a wide range of frequencies. Our own earlier work (White
\& Dey, in prep.; archived at
\url{https://github.com/Tilanthi/SETI}) reproduced and extended the
\citet{Mason2024} methodology directly on ALMA science-target (rather than
calibrator-bycatch) fields for TRAPPIST-1 across Bands 3, 6 and 7, finding no
detection and placing limits of EIRP$_{\rm min}\sim5$--$10\times10^{13}$\,W, and
subsequently built a ranked target list and pilot survey of 100 nearby,
ALMA-observed stars, from which the present, volume-limited survey has grown.
Beyond these efforts, we are not aware of any other published, systematic
technosignature search using ALMA or any comparably sensitive
millimetre/submillimetre interferometer, representing a substantial and readily
addressable gap in frequency coverage relative to the sub-10\,GHz literature.

The near-complete absence of mm/submm technosignature work to date is
attributable less to any fundamental unsuitability of these frequencies -- Sheikh's
``nine axes of merit'' framework \citep{Sheikh2020} notes that higher frequencies
suffer less from interstellar scintillation and dispersion, and offer access to
much larger instantaneous bandwidths -- than to historical circumstance:
sensitive mm/submm interferometers with the necessary spectral resolution and
archival depth (chiefly ALMA) have only existed, and only accumulated a
sufficiently large calibrated science archive, within roughly the last decade.
This makes archival reprocessing, rather than dedicated proposal-driven
observing, a particularly natural strategy in this regime: ALMA's archive already
contains many thousands of hours of high-spectral-resolution observations of
individual, nearby stars, taken for entirely unrelated astrophysical purposes
(disc structure, molecular chemistry, stellar activity, debris-disc detection,
calibration monitoring), each of which can in principle be re-searched for
technosignatures at negligible additional observing cost.

\subsection{Machine learning and anomaly detection in technosignature searches}
\label{sec:ml}

The scale of modern technosignature datasets -- hundreds of hours of
high-time- and frequency-resolution data per survey, often containing many
thousands of candidate ``hits'' dominated by anthropogenic RFI -- has made
machine learning an increasingly central tool for both signal detection and RFI
rejection. \citet{Ma2023} presented the first large-scale deep-learning
technosignature search, applying a $\beta$-convolutional variational autoencoder
to 480 hours of Breakthrough Listen GBT data on 820 nearby stars in a
semi-unsupervised anomaly-detection mode, recovering candidate signals missed by
the standard \textsc{turboSETI} narrowband drift-search pipeline
\citep{Enriquez2019} and reducing the number of signals requiring human vetting by
roughly two orders of magnitude, while returning eight signals of interest for
re-observation. Building on this, \citet{Choza2024} identified and corrected a
long-standing sensitivity mischaracterisation in \textsc{turboSETI}'s
signal-to-noise calibration, prompting a retrospective correction of quoted
sensitivities in several earlier Breakthrough Listen publications
\citep{GBTarchiveML2025}. More recently, unsupervised clustering methods (e.g.\
the ``GLOBULAR'' algorithm) have been developed specifically to separate genuine
anomalies from RFI in narrowband technosignature searches without requiring
labelled training examples \citep{GLOBULAR2025}, and further anomaly-detection
searches have been applied to combined Parkes and GBT Breakthrough Listen archives
\citep{AnomalyParkesGBT2025}. These efforts collectively demonstrate that
machine-learning-based anomaly detection can recover classes of signal
(broadband, non-linearly drifting, or otherwise morphologically unusual) that
narrowband, linear-drift matched-filter searches such as \textsc{turboSETI} are,
by design, insensitive to \citep{BRaTs2026}, and that a purely classical,
threshold-based search -- of the kind we adopt for the present survey, following
\citet{Mason2024} -- necessarily leaves some fraction of the true technosignature
parameter space unexamined. We return to this point in \S\ref{sec:discussion} as
a natural direction for follow-up work once the present survey's baseline dataset
exists.

\subsection{Statistical frameworks for interpreting non-detections}
\label{sec:stats}

A non-detection is not, by itself, a null result: it must be translated into a
quantitative statement about which regions of transmitter parameter space have
been excluded, and with what confidence. \citet{Wright2018} formalised this with
the ``Cosmic Haystack'' framework, modelling the technosignature search space as
an eight-dimensional volume (spanning, e.g., transmitter frequency, bandwidth,
location, distance, and duty cycle) and showing that even the combined efforts of
every major SETI programme conducted up to 2018 sampled a small fraction
of any astrophysically plausible haystack volume -- a result frequently
paraphrased as showing that humanity has searched ``the equivalent of a large
hot tub's worth of the Earth's oceans''. \citet{Sheikh2020} subsequently proposed
a complementary ``nine axes of merit'' rubric for comparing the relative
scientific value of different technosignature search strategies (spanning axes
such as detectability, transmitter cost and information content), providing a
qualitative framework alongside Wright et al.'s quantitative one.
\citet{Margot2023} developed an explicit formalism for converting a survey's
non-detection rate into an upper bound on the fraction of stars in the search
volume hosting a detectable transmitter, but noted explicitly that this
inference is only valid if the observed sample is statistically representative
of the underlying stellar population of the search volume -- a requirement that
most targeted, interest-selected samples (including our own earlier 100-star
pilot list) do not, by design, satisfy. Most recently, \citet{EarthDetectingEarth2025}
inverted the usual question, asking at what distance Earth's own technosignatures
would themselves be detectable by an ALMA- or SKA-like observatory, providing a
valuable calibration point (an ``ichnoscale'' of unity) against which the
sensitivity of any given survey, including the present one, can be benchmarked.

\subsection{The case for a volume-limited, representative sample}
\label{sec:volumelimited}

The great majority of the surveys reviewed above -- including our own earlier
100-star ALMA pilot programme -- select targets by some combination of proximity,
Sun-like spectral type, known exoplanets, protoplanetary-disc or debris-disc
detection, or membership of specially-constructed target lists such as the Earth
Transit Zone. Such selections are entirely reasonable when the goal is to
maximise the a priori probability of detecting a genuine technosignature per
unit observing time, but they carry an unavoidable cost: they make it difficult
to convert a survey's results into a statistically clean statement about the
prevalence of technosignatures in the general stellar population, precisely
because the sample is not representative of that population
\citep{Margot2023}. A disc-hosting or exoplanet-hosting bias is particularly
worth flagging in the ALMA context: ALMA-detectable circumstellar and
protoplanetary discs are, almost without exception, orders of magnitude brighter
than the present-day Solar System's own zodiacal dust cloud would appear at
interstellar distances, so that samples selected \emph{because} they have an
ALMA-detected disc are biased towards young, dynamically active systems that are
not necessarily more (and could plausibly be less) likely to host a mature,
Earth-like technological civilisation than an old, quiescent, disc-poor star of
the same spectral type.

The present survey is explicitly designed to avoid this selection effect. We
construct our sample by identifying every star within 20\,pc of the Sun that
happens to fall within the field of view of at least one public ALMA
observation -- irrespective of whether that observation was proposed to study a
disc, a stellar flare, an unrelated background source, or anything else -- and
irrespective of the target star's own spectral type, age, multiplicity, or known
planet-hosting status. The star need not be at the pointing centre of the
observation, provided it lies within the usable primary-beam field of view (see
\S\ref{sec:sample}). This is, in effect, a volume-limited ``bycatch'' census of
the ALMA archive: rather than sampling stars selected for astrophysical interest
in the calibrator fields as in \citet{Mason2024}, we sample every nearby star
that ALMA happens to have observed for any purpose whatsoever, within a fixed
physical volume. We recognise, and quantify explicitly in \S\ref{sec:sample}
(finding that our 120-star sample represents only $\sim$5\% of the known
stellar population within 20\,pc), that this approach is still subject to
a residual, second-order selection effect -- ALMA archival coverage of the
immediate solar neighbourhood is itself not spatially uniform, since some
regions and spectral types are more likely to have been proposed for ALMA
time than others -- but this bias operates on the presence of \emph{any}
ALMA coverage of a star, not on the astrophysical properties that are the
actual subject of a technosignature search, and can therefore be
characterised and reported as a well-defined completeness function of the
20\,pc volume, rather than an implicit and unquantified bias in the target
selection itself.

\subsection{Prospects for next-generation facilities}
\label{sec:future}

The coming decade will see a substantial expansion of technosignature search
capability. The Square Kilometre Array \citep[SKA;][]{SKA2015,SKA2026review},
now under construction in Australia (SKA-Low) and South Africa (SKA-Mid), will
combine wide instantaneous bandwidth, high angular resolution, and an
architecture well suited to commensal, multi-beam technosignature searching,
with projected sensitivity sufficient to detect Earth-analogue-level radar and
communication leakage from significantly greater distances than is possible
today. The next-generation Very Large Array \citep[ngVLA;][]{ngVLASETI2022} is
specifically being designed with technosignature science as one of its
commensal use cases, with proposed antenna configurations optimised to allow of
order 64 simultaneous high-spectral-resolution commensal beams. Existing
commensal systems at MeerKAT, the VLA (COSMIC), and the Murchison Widefield
Array are already providing a preview of the observing modes these larger
facilities will employ at scale. None of these facilities, however, offer ALMA's
combination of angular resolution, spectral resolution and raw collecting area
above $\sim$\,50\,GHz in the near term, so a systematic mm/submm technosignature
census -- of the kind begun here -- is likely to remain a uniquely ALMA-enabled
measurement for some years yet, and represents a natural complement to the
lower-frequency, larger-sample surveys that will be enabled by SKA and ngVLA.

\section{Sample selection}
\label{sec:sample}

Our sample comprises every star within 20\,pc of the Sun (based on Gaia~DR3
parallaxes) with at least one public ALMA observation, of any band, science
category, or proposal type, in which the star's true (proper-motion-corrected)
position at the epoch of observation falls within the field of view of at least
one execution block -- not necessarily at the observation's pointing centre,
since a substantial fraction of ALMA mosaics and wide single-pointing
observations serendipitously cover more than one catalogued star. We deliberately
do not weight target selection by the presence of a circumstellar disc, known
exoplanets, spectral type, or any other astrophysical property beyond distance
and archival data availability; disc presence and planet-hosting status are
retained purely as descriptive metadata for the final sample, not as
selection or ranking criteria. At the time of writing this identifies 120 unique
target stars, spanning spectral types from early A (e.g.\ Sirius~B, $\alpha$~CMa~B)
through to the latest, faintest M~dwarfs (e.g.\ Proxima Centauri, TRAPPIST-1,
Barnard's Star, Teegarden's Star), ranked by increasing distance from the Sun so
that the nearest, most sensitively-constrained systems are processed first. The
underlying pipeline is built to extend this volume-limited approach to 30, 40 or
50\,pc as a natural next step, at the cost of a rapidly increasing target list
and a correspondingly reduced fraction with existing ALMA coverage, since the
archive's areal completeness at fixed exposure declines with distance.

It is important to be explicit about what ``volume-limited'' does and does
not mean here: our 120 targets are every star within 20\,pc with public
ALMA coverage, \emph{not} a complete census of every known star within
20\,pc. Cross-matching against our own reference Gaia~DR3 catalogue of
the solar neighbourhood, 2357 individual stars are catalogued within
20\,pc, of which our 120 targets represent only $\sim$5\%
(Figure~\ref{fig:sample}). The sample is volume-limited, and unbiased
with respect to spectral type, disc-hosting status, and planet-host
status, \emph{within the population that ALMA happens to have observed
for any purpose}; it is not, and does not claim to be, an unbiased sample
of the local stellar population as a whole, since the decision to propose
ALMA time towards a given nearby star is itself not uniformly
distributed across spectral type or other stellar properties. We regard
this as a genuine, quantifiable limitation of an archival-mining
approach, worth stating plainly rather than allowing the term
``volume-limited'' to be read as ``complete''.

\section{Data and methodology}
\label{sec:method}

For each target we mine the public ALMA science archive for all calibrated (or,
where necessary, locally recalibrated from raw data) execution blocks covering
the star's position, apply an efficiency-bounded download and processing budget
per target to remain within the resource constraints of the shared compute
facility used for this work, and search the resulting visibility data in two
complementary ways: (i) a narrowband, Doppler-drift-corrected search for
technosignature candidates, following the general methodology of
\citet{Mason2024} and \citet{Enriquez2017}, converting any non-detection into an
equivalent isotropic radiated power (EIRP) limit via
$\mathrm{EIRP_{min}} = 4\pi d^{2} S_{\rm min} \Delta\nu$; and (ii) a
line-excluded continuum search for anomalous broad-band emission at the star's
position, with an accompanying epoch-to-epoch variability check wherever
multiple execution blocks exist. Stellar positions are propagated to the epoch
of each observation using Gaia~DR3 proper motions and parallaxes. The
Doppler-drift search range applied to each target is set to the larger of (a) a
generic frequency-scaled default following \citet{Mason2024}'s adoption of a
literature fiducial with a threefold safety margin, and (b) where the target
hosts one or more confirmed exoplanets, an explicit calculation of the maximum
line-of-sight orbital acceleration (and hence maximum plausible transmitter
drift rate) of the innermost known planet, so that a transmitter co-located with
a known, short-period planet cannot be missed simply because its orbital motion
drifts it out of an insufficiently wide search window.

It is important to be explicit about the fraction of each observation's
instantaneous bandwidth that the narrowband search actually covers.
ALMA's correlator is typically configured with several spectral windows
observed simultaneously (a mix of coarse, wide-bandwidth continuum
windows and one or more finer-resolution windows, or, for spectral-scan
observations, many narrow windows tuned across a wide frequency range).
Following \citet{Mason2024}, we search only the single finest-resolution
spectral window available for each target, since narrowband sensitivity
scales with channel width; we do not currently search the coarser
spectral windows configured alongside it. Across the 19 target/bands
searched to date, the number of spectral windows configured
simultaneously ranges from 4 to 24 (median 4), so that the single window
we search represents, on average, only $\sim$23\% (range 4--25\%) of the
instantaneously available spectral windows for a given target -- a
concrete, quantifiable statement of the narrowband search's actual
frequency coverage, which we report per target in
Table~\ref{tab:full} (Appendix~\ref{app:table}) rather than assert only
in aggregate. The line-excluded continuum search (ii), by contrast, does
use the full available bandwidth across all configured spectral windows
for each target, since it does not require narrowband spectral
resolution.

\subsection{Reference transmitter benchmarks}
\label{sec:benchmarks}

To give the EIRP limits reported in \S\ref{sec:results} physical context,
we compare them throughout to two benchmark values for real, existing
human transmitters, both taken from \citet{EarthDetectingEarth2025}: an
Arecibo-like planetary radar, whose characteristic S-band EIRP was
$2\times10^{13}$\,W (the strongest transmitter ever built by humans,
used historically both for Solar System radar astronomy and for the 1974
Arecibo Message); and typical, entirely unintentional Earth radio
leakage from mobile/cellular infrastructure, with an EIRP of order
$4\times10^{9}$\,W. These two values bracket a physically motivated
range: the former represents what a civilisation with our own most
powerful \emph{deliberate, directed} transmitter could achieve if
pointed at us; the latter represents what an entirely unintentional,
Earth-2024-equivalent level of technological activity would radiate in
our direction as background leakage, with no deliberate signalling
intent at all.

\subsection{Closure-phase point-source vetting}
\label{sec:closurephase}

Any narrowband signal that crosses our detection threshold at the star's
position is subjected to a further, automated test that is, to our
knowledge, unavailable to any single-dish or phased-array-beamformed
technosignature survey (Breakthrough Listen, FAST, the ATA, LOFAR, and
similar all lack it, since it requires multiple independent baselines):
the interferometric closure phase, the phase sum
$\arg(V_{ij}) + \arg(V_{jk}) - \arg(V_{ik})$ around a closed triangle of
baselines $i$-$j$-$k$. This quantity is a standard radio-interferometric
identity that is theoretically \emph{exactly zero} for an unresolved
point source at \emph{any} sky position -- a point source's visibility
phase is a pure linear function of $(u,v)$, which cancels identically
around any closed baseline loop -- and is, by construction, immune to
per-station calibration errors. A candidate signal whose closure phase is
statistically inconsistent with zero at the drift-tracked frequency of
the hit is therefore evidence against it behaving as a genuine celestial
point source, favouring instead a near-field, terrestrial, or
instrumental origin. This test is independent of, and complementary to,
the existing spatial-coincidence test (whether the on-source position
shows a larger peak than the control positions): closure phase does not
depend on where the source is, so it answers a different question --
``does this behave like a coherent point source at all'' -- than the
spatial test's ``is it located at the star''. We validated the closure-
phase calculation against synthetic visibility data before applying it to
any ALMA measurement, confirming zero closure phase for a noiseless
simulated point source, statistical consistency with zero for a noisy
simulated point source, and a large, clearly inconsistent scatter for a
simulated non-point-source (RFI-like) signal. Because it is only
triggered on an actual threshold-crossing candidate, this test adds
negligible computational cost to the survey as a whole.

\subsection{Population-level, spectral-type-normalised continuum anomaly check}
\label{sec:popanomaly}

Because our sample is volume-limited and not selected on disc-hosting or
planet-host status, it supports a comparison that is not meaningful for
the biased target lists used in essentially every prior technosignature
survey: whether a star's measured continuum brightness is anomalous
relative to the physically-expected flux of its own stellar photosphere,
and relative to other stars of similar effective temperature already
processed within the same survey. For each target we compute the
expected photospheric flux density from the Planck function and the
star's own solid angle,
$S_\nu = \Omega \times 2h\nu^3/c^2 \times \left[\exp(h\nu/k_BT_{\rm eff}) - 1\right]^{-1}$
with $\Omega = \pi R_\star^2/d^2$, using Gaia DR3 photometric
temperatures and, where available, Gaia GSP-Phot radii (supplemented by a
coarse Pecaut \& Mamajek main-sequence $T_{\rm eff}$--radius relation
where a direct Gaia radius estimate is unavailable, as is the case for
close to half of our sample -- Gaia's astrophysical-parameter pipeline
has known gaps for the very bright/nearby and very red/faint stars that
a 20\,pc sample is disproportionately composed of). Any large excess of
the measured continuum flux (or, for a non-detection, its $5\sigma$ upper
limit) above this photospheric baseline is flagged, and is additionally
compared against the distribution of excess ratios already measured for
other stars in the same coarse $T_{\rm eff}$ bin, once at least five such
comparison stars exist, using a robust (median/MAD-based) outlier
statistic. This check requires no additional raw-data retention beyond
what the continuum measurement itself already produces, and updates
incrementally as the survey proceeds.

\section{Results}
\label{sec:results}

As of this version of the paper, 19 of the planned 120 target/bands have
been searched to completion (a further 2 completed their continuum
measurement but not the narrowband search, owing to a data-quality
failure in those specific execution blocks; both are flagged for an
automatic, non-gating retry). The survey is strictly volume-ordered
(nearest star first; \S\ref{sec:sample}), so the results reported here
are weighted towards the very nearest stars in the sample, spanning
1.3--19.6\,pc at the time of writing. Full per-target values, including
the exact frequency range searched in each band, are given in
Table~\ref{tab:full} in Appendix~\ref{app:table} rather than in the main
text; this section summarises the results and highlights points of
interest.

\subsection{Technosignature search}

Seventeen execution blocks yielded a narrowband search result (two
further targets, Barnard's Star and Wolf~359, had their continuum
measurement succeed but their narrowband search fail outright, and carry
no EIRP limit at this time). The resulting EIRP$_{\rm min}$ limits span
$1.6\times10^{13}$--$2.8\times10^{15}$\,W, with a median of
$2.1\times10^{14}$\,W (Figure~\ref{fig:eirp}, left panel); as expected
for a volume-ordered survey, the deepest limits are for the very nearest
systems (e.g.\ $1.7\times10^{13}$\,W for the UV~Ceti binary at $\sim$2.7\,pc).
Compared against the reference transmitter benchmarks of
\S\ref{sec:benchmarks}, our very deepest limits (the handful of targets
within $\sim$3\,pc) are only just below the $2\times10^{13}$\,W
Arecibo-like planetary-radar benchmark, meaning we could have detected an
equally powerful, deliberately Earth-directed transmitter at those
distances had one been present; every target beyond $\sim$4\,pc, and
every target with respect to the much fainter $\sim4\times10^{9}$\,W
level of Earth's own unintentional radio leakage, remains well out of
reach of the present survey's sensitivity (Figure~\ref{fig:eirp}). This
is a useful, sobering calibration: even this deep, targeted, archival
mm-wave search is not yet sensitive to an Earth-equivalent civilisation's
background leakage at any of the distances probed so far, and is only
marginally sensitive to an Earth-equivalent civilisation's most powerful
deliberate transmitter, and only for its very nearest neighbours.
One narrowband hit was automatically flagged by the pipeline as a
``credible technosignature candidate'' (on-source peak signal-to-noise
ratio above threshold and localised to the star's position, rather than
appearing equally in the control positions): a signal towards $\beta$~Pictoris
in ALMA Band~3 at 115.26\,GHz. This frequency coincides, to within the
spectral resolution of the observation, with the rest frequency of the
CO($J=1\to0$) rotational transition (115.271\,GHz), and $\beta$~Pictoris
is a well-known debris-disc system with previously reported CO emission;
we therefore attribute this hit to line contamination rather than a
genuine technosignature. No other candidate has survived both the
spatial-coincidence and (where triggered) closure-phase vetting tests
described in \S\ref{sec:closurephase}; in practice, the closure-phase test
has not yet been exercised against a real ALMA measurement, since no
other candidate has crossed the initial detection threshold to trigger it.

\subsection{Continuum search}

Of the 19 target/bands with a continuum measurement, five show a
$\geq5\sigma$ continuum detection and fourteen are non-detections
(Figure~\ref{fig:eirp}, right panel). Throughout this paper, a continuum
``upper limit'' is explicitly $5\times$ the map rms noise at the star's
position (i.e.\ the same $5\sigma$ threshold used to define a detection
in the first place); the underlying rms values themselves are also given
per target in Table~\ref{tab:full} so that a reader wanting a different
confidence level can rescale directly, rather than being limited to a
single quoted multiple of $\sigma$. On this basis the non-detections
reach 5$\sigma$ upper limits of 0.19--3.6\,mJy (median 0.59\,mJy). Two of the five
detections -- catalogued separately in our target list as G~272-61A and
G~272-61B -- are in fact the two components of the well-known, closely
separated visual binary UV~Ceti, observed within a single ALMA pointing;
the reported flux (0.85\,mJy) is for the spatially unresolved, blended
pair and cannot be attributed to either component individually. The
remaining three detections ($\tau$~Ceti, and $\beta$~Pictoris in both
Bands 3 and 6) are consistent with previously known circumstellar/debris
material or stellar activity in these well-studied systems and are not
treated as anomalous. The population-level anomaly check described in
\S\ref{sec:popanomaly} has been applied to every target processed to
date but has not yet flagged any star, and in most $T_{\rm eff}$ bins the
survey does not yet have the minimum of five comparison stars needed for
the population-relative component of the test; both the physical-baseline
and population-relative statistics will accumulate statistical power as
the survey proceeds.

\begin{figure*}
\centering
\includegraphics[width=0.95\textwidth]{figures/eirp_continuum_vs_distance.pdf}
\caption{Technosignature EIRP upper limits (left) and continuum flux
densities (right; downward triangles are $5\sigma$ upper limits, i.e.\
5$\times$ the map rms noise -- see \S\ref{sec:results} and
Table~\ref{tab:full} for the rms values themselves -- filled circles are
detections) as a function of distance, for the 19 target/bands searched
to date. The single automatically-flagged candidate ($\beta$~Pictoris,
Band 3) is marked with a star and is identified in the text as CO(1--0)
line contamination rather than a technosignature. The left panel also
shows two reference transmitter benchmarks from
\citet{EarthDetectingEarth2025} (\S\ref{sec:benchmarks}): the EIRP of an
Arecibo-like planetary radar ($2\times10^{13}$\,W, the most powerful
deliberate transmitter yet built by humans) and of typical, unintentional
Earth radio leakage ($\sim4\times10^{9}$\,W).}
\label{fig:eirp}
\end{figure*}

\begin{figure}
\centering
\includegraphics[width=\columnwidth]{figures/sample_distance_distribution.pdf}
\caption{Cumulative number of stars, $N(<d)$, within 20\,pc as a function
of distance $d$. The grey curve is the full known stellar population
within 20\,pc from our reference Gaia~DR3 catalogue (2357 stars); the
blue curve is our 120-star volume-limited sample -- every one of those
2357 stars that additionally has public ALMA archival coverage, only
$\sim$5\% of the total. The two curves are not expected to track one
another: ALMA proposal pressure on the immediate solar neighbourhood is
itself not spectral-type- or distance-uniform, so our sample is
volume-limited and unbiased \emph{within} the ALMA-observed population,
but is not a complete or unbiased sample of the local stellar population
as a whole (see \S\ref{sec:sample}).}
\label{fig:sample}
\end{figure}

\section{Discussion}
\label{sec:discussion}

The results presented here are necessarily preliminary: only 19 of 120
planned target/bands (16\%) have been completed, and the survey is
ordered to complete the nearest, most sensitively-constrained systems
first, so the eventual full sample will extend these results to
significantly larger distances and, correspondingly, shallower EIRP and
continuum limits at the faint end. No credible technosignature has been
found among the systems searched so far; the one automatically-flagged
candidate is confidently attributable to CO(1--0) line emission rather
than an artificial signal, underlining the importance -- already
highlighted by \citet{Mason2024} -- of checking any millimetre-band
narrowband hit against known molecular transitions before treating it as
a candidate. Applying the statistical frameworks of \citet{Wright2018}
and \citet{Margot2023} to derive a quantitative bound on the prevalence
of detectable transmitters in the solar neighbourhood is deferred to a
future version of this paper, once a statistically meaningful fraction
of the volume-limited sample has been completed; with only 19 systems
searched to date such a bound would not yet be usefully constraining.

The two supplementary analyses introduced in this work
(\S\ref{sec:closurephase}, \S\ref{sec:popanomaly}) are best understood as
infrastructure investments rather than sources of results in their own
right at this early stage. The closure-phase test has not yet been
exercised against real data, simply because no candidate other than the
already-explained $\beta$~Pictoris CO line has crossed threshold; it
remains available, at negligible ongoing cost, for any future hit. The
population-level anomaly check is accumulating the population of
same-$T_{\rm eff}$-bin comparison stars it needs to become statistically
powerful, and we expect it to become genuinely diagnostic only once
several dozen targets per broad spectral-type bin have been processed.
We regard both as a useful methodological contribution independent of
whether they ultimately flag anything in this particular sample: to our
knowledge, neither has previously been applied in a technosignature
search, and both are made possible specifically by properties of this
survey (interferometric data; a volume-limited, unbiased sample) that do
not hold for the great majority of prior technosignature work reviewed in
\S\ref{sec:background}.

The UV~Ceti unresolved-pair case (\S\ref{sec:results}) illustrates a
general caveat for any archival, position-agnostic survey of this kind:
close visual binaries and multiple systems that are unresolved at the
relevant ALMA angular resolution will necessarily share a single
measurement, and any future statistical analysis of this sample's
completeness must account for this rather than double-counting such
pairs as independent trials.

\section{Conclusions}
\label{sec:conclusions}

We have described the design, methodology, and first results of a
volume-limited technosignature search of all 120 stars within 20\,pc
with public ALMA archival coverage, irrespective of disc-hosting or
planet-host status. Nineteen of the planned 120 target/bands have been
completed as of this version, yielding EIRP limits of
$1.6\times10^{13}$--$2.8\times10^{15}$\,W and fourteen continuum
non-detections reaching a median $5\sigma$ limit of 0.59\,mJy, with no
credible technosignature identified. We have introduced two analyses
-- interferometric closure-phase vetting and a population-level,
spectral-type-normalised continuum anomaly check -- that exploit
properties unique to this survey and are, to our knowledge, new to the
technosignature literature; both are now running as a standard part of
the pipeline and will accumulate statistical power as the survey
proceeds. This paper will be updated as further targets are completed,
and we anticipate extending the volume limit to 30, 40, and 50\,pc as a
natural next step once the 20\,pc sample is complete.

\section*{Acknowledgements}
This paper makes use of the following ALMA data: [the full list of ALMA
project codes used will be compiled once the survey is complete; see
Table~\ref{tab:full} in Appendix~\ref{app:table} for the individual
execution blocks searched to date]. ALMA is a partnership of ESO
(representing its member states), NSF (USA) and NINS (Japan), together with NRC
(Canada), MOST and ASIAA (Taiwan), and KASI (Republic of Korea), in cooperation
with the Republic of Chile. The Joint ALMA Observatory is operated by ESO,
AUI/NRAO and NAOJ. This work has made use of data from the European Space
Agency (ESA) mission {\it Gaia}, processed by the {\it Gaia} Data Processing and
Analysis Consortium (DPAC).

\section*{Data Availability}
The reduced data products underlying this paper, together with the full
target-selection and analysis pipeline, are being made available at
\url{https://github.com/Tilanthi/SETI} as the survey progresses. Raw ALMA
visibility data are publicly available from the ALMA Science Archive at
\url{https://almascience.org}.


\appendix
\section{Full per-target results table}
\label{app:table}

Table~\ref{tab:full} lists the full per-target/band results underlying the
summary given in Section~\ref{sec:results}: the narrowband technosignature
frequency range actually searched, the number of spectral windows
configured simultaneously by the correlator for that observation
(``N$_{\rm spw}$''; only one of these -- the finest-resolution window --
is ever searched narrowband, per \S\ref{sec:method}), the resulting
EIRP$_{\rm min}$ limit, whether the automated pipeline flagged a credible
candidate, the continuum image rms noise, and the continuum flux density
or, for non-detections, its $5\sigma$ (i.e.\ $5\times$rms) upper limit,
denoted with the standard ``$<$'' notation. This table will grow as the
survey proceeds; the version current at the time of writing is
reproduced here so that the summary statistics quoted in the main text
are fully traceable.

\begin{table*}
\centering
\scriptsize
\caption{Per-target/band results, ordered by distance. Frequency range is
in GHz; N$_{\rm spw}$ is the total number of spectral windows configured
simultaneously in the observation (only 1 of these is ever searched
narrowband; \S\ref{sec:method}); EIRP$_{\rm min}$ is in W. ``Cand.''
indicates whether the automated pipeline flagged the narrowband search as
a credible technosignature candidate (Y/N); see \S\ref{sec:results} for
the disposition of the one Y in this table. Continuum values are the map
rms noise and either the measured peak flux density (for a detection) or
the resulting $5\sigma$ upper limit (denoted ``$<$''), both in mJy.}
\label{tab:full}
\begin{tabularx}{\textwidth}{@{}l c c c c c c c c X@{}}
\hline
Target & Band & Dist.\ (pc) & Freq.\ range (GHz) & N$_{\rm spw}$ & EIRP$_{\rm min}$ (W) & Cand.\ & RMS (mJy) & Flux (mJy) & Notes \\
\hline
Proxima Cen & B6 & 1.30 & 223.09--224.91 & 4 & 3.43e+13 & N & 0.1332 & $<$0.666 &  \\
NAME Barnard's star & B7 & 1.83 & \ldots & 6 & \ldots & \ldots & 0.7243 & $<$3.622 & narrowband search failed \\
Wolf  359 & B7 & 2.41 & \ldots & 7 & \ldots & \ldots & 0.6761 & $<$3.380 & narrowband search failed \\
Sirius B & B3 & 2.67 & 92.13--93.86 & 4 & 2.56e+13 & N & 0.0379 & $<$0.189 &  \\
Sirius B & B4 & 2.67 & 142.13--143.86 & 4 & 2.31e+13 & N & 0.0493 & $<$0.247 &  \\
Sirius B & B5 & 2.67 & 195.13--196.86 & 4 & 3.48e+13 & N & 0.0911 & $<$0.455 &  \\
G 272-61B & B3 & 2.67 & 89.63--91.36 & 4 & 1.65e+13 & N & 0.0260 & 0.850 & unresolved UV Cet AB pair, blended flux \\
G 272-61A & B3 & 2.72 & 89.63--91.36 & 4 & 1.71e+13 & N & 0.0260 & 0.850 & unresolved UV Cet AB pair, blended flux \\
CD-23 14742 & B6 & 2.98 & 223.09--224.91 & 4 & 1.20e+14 & N & 0.0776 & $<$0.388 &  \\
tau Cet &  & 3.65 & 229.63--231.47 & 24 & 3.79e+13 & N & 0.1066 & 0.863 &  \\
GJ 674 & B6 & 4.55 & 223.09--224.91 & 4 & 2.74e+14 & N & 0.1032 & $<$0.516 &  \\
GJ 581 &  & 6.30 & 223.09--224.91 & 4 & 6.35e+14 & N & 0.0903 & $<$0.451 &  \\
GJ 849 &  & 8.81 & 229.08--230.92 & 4 & 8.30e+14 & N & 0.1424 & $<$0.712 &  \\
HD 285968 &  & 9.49 & 229.04--230.88 & 4 & 3.45e+14 & N & 0.1400 & $<$0.700 &  \\
HD 23484 & B6 & 16.17 & 229.61--231.45 & 4 & 1.70e+15 & N & 0.0995 & $<$0.497 &  \\
HD 10647 & B7 & 17.35 & 330.35--332.19 & 4 & 1.97e+15 & N & 0.4040 & $<$2.020 &  \\
HD 53143 &  & 18.34 & 229.60--231.44 & 4 & 2.84e+15 & N & 0.1754 & $<$0.877 &  \\
bet Pic & B3 & 19.63 & 114.83--115.69 & 4 & 2.08e+14 & Y & 0.0281 & 0.444 & CO(1-0) line, see text \\
bet Pic & B6 & 19.63 & 230.08--231.00 & 4 & 2.01e+15 & N & 0.9626 & 5.659 &  \\
\hline
\end{tabularx}
\end{table*}

\begin{thebibliography}{99}

\bibitem[Backus \& Project Phoenix Team(2002)]{Backus2002} Backus P.~R., Project Phoenix Team, 2002, in Bulletin of the American Astronomical Society. p.\ 1173

\bibitem[Choza et al.(2024)]{Choza2024} Choza C., et al., 2024, AJ, 168, 254

\bibitem[Cocconi \& Morrison(1959)]{CocconiMorrison1959} Cocconi G., Morrison P., 1959, Nature, 184, 844

\bibitem[Czech et al.(2021)]{Czech2021} Czech D., et al., 2021, PASP, 133, 064502

\bibitem[Czech et al.(2026)]{BLMeerKAT2026} Czech D.~J., et al., 2026, arXiv:2607.23651

\bibitem[Drake(1960)]{Drake1960} Drake F.~D., 1960, Physics Today, 13, 40

\bibitem[Drake(1974)]{Drake1974} Drake F.~D., 1974, in Sagan C., ed., Communication with Extraterrestrial Intelligence (CETI). MIT Press

\bibitem[Drake(1986)]{Drake1986} Drake F.~D., 1986, in Applications of Digital Image Processing IX. SPIE

\bibitem[Enriquez \& Price(2019)]{Enriquez2019} Enriquez E., Price D., 2019, turboSETI, Astrophysics Source Code Library, record ascl:1906.006

\bibitem[Enriquez et al.(2017)]{Enriquez2017} Enriquez J.~E., et al., 2017, ApJ, 849, 104

\bibitem[Gajjar et al.(2019)]{Gajjar2019EarthTransitZone} Gajjar V., et al., 2019, in Bulletin of the American Astronomical Society. AAS Meeting \#233

\bibitem[Garrett et al.(2026)]{BRaTs2026} Garrett M.~A., et al., 2026, MNRAS, accepted (arXiv:2605.10212)

\bibitem[Isaacson et al.(2017)]{BLoverview} Isaacson H., et al., 2017, PASP, 129, 054501

\bibitem[Jacobson-Bell et al.(2025)]{GLOBULAR2025} Jacobson-Bell K., et al., 2025, AJ, 169, 233

\bibitem[Johnson et al.(2023)]{LOFARdualsite2023} Johnson O.~A., et al., 2023, AJ, 166, 193

\bibitem[Lacki et al.(2021)]{GBTarchiveML2025} Lacki B.~C., et al., 2021, arXiv:2101.02244

\bibitem[Li et al.(2026)]{FAST3IATLAS2025} Li J.-K., Tao Z.-Z., Zhang T.-J., 2026, arXiv:2603.19023

\bibitem[Ma et al.(2023)]{Ma2023} Ma P.~X., et al., 2023, Nature Astronomy, 7, 492

\bibitem[Margot et al.(2023)]{Margot2023} Margot J.-L., et al., 2023, AJ, 166, 189

\bibitem[Mason et al.(2024)]{Mason2024} Mason L.~A., Garrett M.~A., Wandia K., Siemion A.~P.~V., 2024, MNRAS, 536, 2127

\bibitem[Morrison et al.(2023)]{Morrison2023} Morrison I.~S., et al., 2023, PASA, 40, e038

\bibitem[Ng et al.(2022)]{ngVLASETI2022} Ng C., et al., 2022, AJ, 164, 205

\bibitem[Poznanski et al.(2025)]{AnomalyParkesGBT2025} Poznanski D., et al., 2025, AJ, in press (arXiv:2505.03927)

\bibitem[Price et al.(2020)]{Price2020} Price D.~C., et al., 2020, AJ, 159, 86

\bibitem[Sheikh(2020)]{Sheikh2020} Sheikh S.~Z., 2020, International Journal of Astrobiology, 19, 237

\bibitem[Sheikh et al.(2025a)]{EarthDetectingEarth2025} Sheikh S.~Z., et al., 2025, AJ, 169, 108

\bibitem[Sheikh et al.(2025b)]{ATAISO2025} Sheikh S.~Z., et al., 2025, arXiv:2512.18142 (``A Search for Radio Technosignatures from Interstellar Object 3I/ATLAS with the Allen Telescope Array'')

\bibitem[Siemion et al.(2015)]{SKA2015} Siemion A.~P.~V., et al., 2015, arXiv:1412.4867

\bibitem[Tarter(2001)]{Tarter2001} Tarter J., 2001, ARA\&A, 39, 511

\bibitem[Tremblay et al.(2024)]{Tremblay2024} Tremblay C.~D., et al., 2024, PASA, 41, e004

\bibitem[Tremblay et al.(2026)]{SKA2026review} Tremblay C.~D., et al., 2026, in Advancing Astrophysics with the SKA II (AASKAII), arXiv:2606.27565

\bibitem[Tusay et al.(2024)]{ATATRAPPIST2024} Tusay N., Sheikh S.~Z., Sneed E.~L., Farah W., Pollak A.~W., Cruz L.~F., Siemion A., DeBoer D.~R., Wright J.~T., 2024, AJ, 168, 283

\bibitem[Tusay et al.(2025)]{TESSeclipsing2025} Tusay N., et al., 2025, arXiv:2506.13459

\bibitem[Valente et al.(2025)]{SRT2025} Valente G., et al., 2025, Acta Astronautica, in press

\bibitem[Vidal et al.(2026)]{SETIreview2025} Vidal C., et al., 2026, arXiv:2605.21093 (``The Search for Technosignatures: a Review of Possibilities'')

\bibitem[Von Korff et al.(2025)]{SETIhome2025} Von Korff J., et al., 2025, AJ, 170, 112

\bibitem[Wright et al.(2018)]{Wright2018} Wright J.~T., Kanodia S., Lubar E.~G., 2018, AJ, 156, 260

\bibitem[Zhang et al.(2020)]{Zhang2020} Zhang Z.-S., et al., 2020, ApJ, 891, 174

\end{thebibliography}

\end{document}
